\BeleskeID{09_2-s029}
% element: 09_2-s029:e04
Daljim povećanjem ulaznog napona izlazni napon opada, dok napon između sorsa i drejna P tranzistora raste. Oblast svakog tranzistora određuje se poređenjem napona drejn–sors sa granicom zasićenja njegovog modela. U oznakama priključaka ti uslovi su
\[V_{DS,N} \ge \frac{(V_{GS,N}-V_{Tn})}{1+\frac{(V_{GS,N}-V_{Tn})}{L_nE_{Cn}}}\]
\[V_{DS,P} \le \frac{(V_{GS,P}-V_{Tp})}{1+\frac{(V_{GS,P}-V_{Tp})}{L_pE_{Cp}}}\]
P tranzistor može preći u zasićenje dok N još ostaje zasićen. U toj oblasti oba izraza za struju koriste model zasićenja.
